% !TeX root = ../colloquium-2.tex

% Определение производной порядка больше или равно 2. 
% Формулировка обобщённого правила Лейбница.

\phantomsection\label{sec_Definition15}
\begin{definition}[
    Определение производной порядка больше или равно 2.
    Формулировка обобщённого правила Лейбница.
]
    \subsection*{Определение производной порядка больше или равно 2.}
    Пусть функция $f$ дифференцируема в каждой точке множества $E$. \\
    Тогда её производная $f'$ определена на $E$ и можно ставить вопрос о её 
    дифференцируемости. \\
    Отметим, что производная функции $f$ нулевого порядка называется сама функция $f$.
    Производная порядка $n$ определяется индуктивно.
    \begin{note}[purple]
        Если у функции $f$ в некоторой окрестности точки $a$ существует производная
        $n - 1$-го порядка $f^{(n - 1)}$ (то есть эта производная представляет собой
        определённую в некоторой окрестности точки $a$ функцию) и эта производная 
        дифференцируема в точке $a$, то производная функции $f$ порядка $n$ в точке $a$ - 
        это производная функции $f^{(n - 1)}$ в точке $a \; : \; f^{(n)}(a) := (f^{(n - 1)})'(a)$. \\

        \textbf{В терминах пределов:}
        $$ 
            f^{(n)}(a) = \lim_{h \rightarrow 0} \frac{f^{(n - 1)}(a + h) - f^{(n - 1)}(a)}{h}
        $$

        Ещё раз отметим, что производная порядка $n - 1$ определена в окрестности точки $a$, 
        иначе символ $f^{(n - 1)} (a + h)$ не имеет смысла.
    \end{note}

    \subsection*{Формулировка обобщённого правила Лейбница.}
    \begin{note}[purple]
        Пусть функции $u$ и $v$ $n$ раз дифференцируемы в точке $a$. 
        Тогда их произведение также $n$ раз дифференцируемо в точке $a$ и
        $$ (u \cdot v)^{(n)} (a) = \sum_{k = 0}^{n} C_n^k u^{(k)} (a) v^{(n - k)}(a) $$
    \end{note}
\end{definition}